\documentclass{amsart}
% For dealing with references we use the comment environment
\usepackage{verbatim}
\newenvironment{reference}{\comment}{\endcomment}
%\newenvironment{reference}{}{}
\newenvironment{slogan}{\comment}{\endcomment}
\newenvironment{history}{\comment}{\endcomment}

% For commutative diagrams we use Xy-pic
\usepackage[all]{xy}

% We use 2cell for 2-commutative diagrams.
\xyoption{2cell}
\UseAllTwocells

% We use multicol for the list of chapters between chapters
\usepackage{multicol}

% This is generally recommended for better output
\usepackage{lmodern}
\usepackage[T1]{fontenc}

% For cross-file-references

% Package for hypertext links:
\usepackage{hyperref}

% For any local file, say "hello.tex" you want to link to please

% Theorem environments.
%
\theoremstyle{plain}
\newtheorem{theorem}[subsection]{Theorem}
\newtheorem{proposition}[subsection]{Proposition}
\newtheorem{lemma}[subsection]{Lemma}

\theoremstyle{definition}
\newtheorem{definition}[subsection]{Definition}
\newtheorem{example}[subsection]{Example}
\newtheorem{exercise}[subsection]{Exercise}
\newtheorem{situation}[subsection]{Situation}

\theoremstyle{remark}
\newtheorem{remark}[subsection]{Remark}
\newtheorem{remarks}[subsection]{Remarks}

\numberwithin{equation}{subsection}

% Macros
%
\def\lim{\mathop{\mathrm{lim}}\nolimits}
\def\colim{\mathop{\mathrm{colim}}\nolimits}
\def\Spec{\mathop{\mathrm{Spec}}}
\def\Hom{\mathop{\mathrm{Hom}}\nolimits}
\def\Ext{\mathop{\mathrm{Ext}}\nolimits}
\def\SheafHom{\mathop{\mathcal{H}\!\mathit{om}}\nolimits}
\def\SheafExt{\mathop{\mathcal{E}\!\mathit{xt}}\nolimits}
\def\Sch{\mathit{Sch}}
\def\Mor{\mathop{\mathrm{Mor}}\nolimits}
\def\Ob{\mathop{\mathrm{Ob}}\nolimits}
\def\Sh{\mathop{\mathit{Sh}}\nolimits}
\def\NL{\mathop{N\!L}\nolimits}
\def\CH{\mathop{\mathrm{CH}}\nolimits}
\def\proetale{{pro\text{-}\acute{e}tale}}
\def\etale{{\acute{e}tale}}
\def\QCoh{\mathit{QCoh}}
\def\Ker{\mathop{\mathrm{Ker}}}
\def\Im{\mathop{\mathrm{Im}}}
\def\Coker{\mathop{\mathrm{Coker}}}
\def\Coim{\mathop{\mathrm{Coim}}}

% Boxtimes
%
\DeclareMathSymbol{\boxtimes}{\mathbin}{AMSa}{"02}

%
% Macros for moduli stacks/spaces
%
\def\QCohstack{\mathcal{QC}\!\mathit{oh}}
\def\Cohstack{\mathcal{C}\!\mathit{oh}}
\def\Spacesstack{\mathcal{S}\!\mathit{paces}}
\def\Quotfunctor{\mathrm{Quot}}
\def\Hilbfunctor{\mathrm{Hilb}}
\def\Curvesstack{\mathcal{C}\!\mathit{urves}}
\def\Polarizedstack{\mathcal{P}\!\mathit{olarized}}
\def\Complexesstack{\mathcal{C}\!\mathit{omplexes}}
% \Pic is the operator that assigns to X its picard group, usage \Pic(X)
% \Picardstack_{X/B} denotes the Picard stack of X over B
% \Picardfunctor_{X/B} denotes the Picard functor of X over B
\def\Pic{\mathop{\mathrm{Pic}}\nolimits}
\def\Picardstack{\mathcal{P}\!\mathit{ic}}
\def\Picardfunctor{\mathrm{Pic}}
\def\Deformationcategory{\mathcal{D}\!\mathit{ef}}

\setlength{\emergencystretch}{3em}
\begin{document}
\title[Stacks Project, Tag 06T8]{401 --- Minimal versal tangent maps are bijective modulo residue-field derivations}
\maketitle
\noindent Copyright (C) 2005--2025 Johan de Jong. Stacks Project authors. Source: \href{https://stacks.math.columbia.edu/tag/06T8}{Tag 06T8}.

\noindent Editorial difficulty note (not author text). Difficulty H2: The proof detects a nontrivial tangent kernel outside the residue-derivation orbits and forms the kernel subring of that derivation. It decomposes the first-order algebra as a dual-number product and uses the Schlessinger lifting condition to descend the object to a proper subring. This contradicts minimality and establishes the orbit-level injectivity even for inseparable residue fields..

\noindent Editorial provenance note (not author text). History: excerpt from revision \texttt{a0444\allowbreak 6e57e\allowbreak c1fbc\allowbreak 252a8\allowbreak 71afc\allowbreak ec775\allowbreak 2fb28\allowbreak 07b14}, formal-defos.tex, lines 4526--4617. Reference presentation was adapted; original-author context and core support blocks were retained or added. The mathematical wording is unchanged. Licensed under GNU FDL 1.2 or later; no invariant sections or cover texts. See ../licenses/COPYING. Context blocks reproduce author definitions and supporting results; they are not additional counted problems.

\begin{definition}
\label{definition-CLambda}
Let $\Lambda$ be a Noetherian ring and let $\Lambda \to k$ be a finite
ring map where $k$ is a field. We define {\it $\mathcal{C}_\Lambda$} to be
the category with
\begin{enumerate}
\item objects are pairs $(A, \varphi)$ where $A$ is an Artinian local
$\Lambda$-algebra and where $\varphi : A/\mathfrak m_A \to k$ is a
$\Lambda$-algebra isomorphism, and
\item morphisms $f : (B, \psi) \to (A, \varphi)$ are local $\Lambda$-algebra
homomorphisms such that $\varphi \circ (f \bmod \mathfrak m) = \psi$.
\end{enumerate}
We say we are in the {\it classical case} if $\Lambda$ is a Noetherian
complete local ring and $k$ is its residue field.
\end{definition}

\section{Notation and Conventions}
\label{section-notations-conventions}

\noindent
A ring is commutative with $1$. The  maximal ideal of a local ring $A$
is denoted by $\mathfrak{m}_A$. The set of positive integers is denoted
by $\mathbf{N} = \{1, 2, 3, \ldots\}$. If $U$ is an object of a
category $\mathcal{C}$, we denote by $\underline{U}$
the functor
$\Mor_\mathcal{C}(U, -): \mathcal{C} \to \textit{Sets}$, see
Remarks \href{https://stacks.math.columbia.edu/tag/06GK}{remarks cofibered groupoids (Tag 06GK)} (\href{https://stacks.math.columbia.edu/tag/06GQ}{item definition yoneda (Tag 06GQ)}).
Warning: this may conflict with the notation in other chapters where we
sometimes use $\underline{U}$ to denote $h_U(-) = \Mor_\mathcal{C}(-, U)$.

\medskip\noindent
Throughout this chapter $\Lambda$ is a Noetherian ring and
$\Lambda \to k$ is a finite ring map from $\Lambda$ to a field.
The kernel of this map is denoted $\mathfrak m_\Lambda$ and the
image $k' \subset k$. It turns out that $\mathfrak m_\Lambda$ is
a maximal ideal, $k' = \Lambda/\mathfrak m_\Lambda$ is a field, and
the extension $k/k'$ is finite. See discussion surrounding
(\href{https://stacks.math.columbia.edu/tag/06S2}{equation k prime (Tag 06S2)}).




\begin{definition}
\label{definition-completion-CLambda}
Let $\Lambda$ be a Noetherian ring and let $\Lambda \to k$ be a finite
ring map where $k$ is a field. We define {\it $\widehat{\mathcal{C}}_\Lambda$}
to be the category with
\begin{enumerate}
\item objects are pairs $(R, \varphi)$ where $R$ is a Noetherian complete
local $\Lambda$-algebra and where $\varphi : R/\mathfrak m_R \to k$ is a
$\Lambda$-algebra isomorphism, and
\item morphisms $f : (S, \psi) \to (R, \varphi)$ are local $\Lambda$-algebra
homomorphisms such that $\varphi \circ (f \bmod \mathfrak m) = \psi$.
\end{enumerate}
\end{definition}

\begin{definition}
\label{definition-category-cofibred-groupoids}
Let $\mathcal{C}$ be a category.  A {\it category cofibered in groupoids over
$\mathcal{C}$} is a category $\mathcal{F}$ equipped with a functor
$p: \mathcal{F} \to \mathcal{C}$ such that $\mathcal{F}^{opp}$ is a category
fibered in groupoids over $\mathcal{C}^{opp}$ via
$p^{opp}: \mathcal{F}^{opp} \to \mathcal{C}^{opp}$.
\end{definition}

\begin{definition}
\label{definition-predeformation-category}
A {\it predeformation category} $\mathcal{F}$ is a category cofibered
in groupoids over $\mathcal{C}_\Lambda$ such that $\mathcal{F}(k)$ is
equivalent to a category with a single object and a single morphism,
i.e., $\mathcal{F}(k)$ contains at least one object and there is a
unique morphism between any two objects. A {\it morphism of predeformation
categories} is a morphism of categories cofibered in groupoids over
$\mathcal{C}_\Lambda$.
\end{definition}

\begin{definition}
\label{definition-formal-objects}
Let $\mathcal{F}$ be a category cofibered in groupoids over
$\mathcal{C}_\Lambda$. The {\it category $\widehat{\mathcal{F}}$ of formal
objects of  $\mathcal{F}$} is the category with the following objects and
morphisms.
\begin{enumerate}
\item A {\it formal object $\xi = (R, \xi_n, f_n)$ of $\mathcal{F}$}
consists of an object $R$ of $\widehat{\mathcal{C}}_\Lambda$, and a collection
indexed by $n \in \mathbf{N}$ of objects $\xi_n$ of
$\mathcal{F}(R/\mathfrak m_R^n)$ and morphisms
$f_n : \xi_{n + 1} \to \xi_n$ lying over the projection
$R/\mathfrak m_R^{n + 1} \to R/\mathfrak m_R^n$.
\item Let $\xi = (R, \xi_n, f_n)$ and $\eta = (S, \eta_n, g_n)$ be
formal objects of $\mathcal{F}$.  A {\it morphism $a : \xi \to \eta$ of
formal objects} consists of a map $a_0 : R \to S$ in
$\widehat{\mathcal{C}}_\Lambda$ and a collection $a_n : \xi_n \to \eta_n$
of morphisms of $\mathcal{F}$ lying over
$R/\mathfrak m_R^n \to S/\mathfrak m_S^n$,
such that for every $n$ the diagram
$$
\xymatrix{
\xi_{n + 1} \ar[r]_{f_n} \ar[d]_{a_{n + 1}} & \xi_n \ar[d]^{a_n} \\
\eta_{n + 1} \ar[r]^{g_n} & \eta_n
}
$$
commutes.
\end{enumerate}
\end{definition}

\begin{definition}
\label{definition-versal}
Let $\mathcal{F}$ be a category cofibered in groupoids.  Let $\xi$ be a formal
object of $\mathcal{F}$ lying over $R \in \Ob(\widehat{\mathcal{C}}_\Lambda)$.
We say $\xi$ is {\it versal} if the corresponding morphism
$\underline{\xi}: \underline{R}|_{\mathcal{C}_\Lambda} \to \mathcal{F}$
of Remark \href{https://stacks.math.columbia.edu/tag/06HC}{remark formal objects yoneda (Tag 06HC)} is smooth.
\end{definition}

\begin{definition}
\label{definition-S1-S2}
Let $\mathcal{F}$ be a category cofibered in groupoids over $\mathcal
C_\Lambda$. We define {\it conditions (S1) and (S2)}
on $\mathcal{F}$ as follows:
\begin{enumerate}
\item[(S1)] Every diagram in $\mathcal{F}$
$$
\vcenter{
\xymatrix{
           & x_2 \ar[d] \\
x_1 \ar[r] & x
}
}
\quad\text{lying over}\quad
\vcenter{
\xymatrix{
           & A_2 \ar[d] \\
A_1 \ar[r] & A
}
}
$$
in $\mathcal{C}_\Lambda$ with $A_2 \to A$ surjective can be completed
to a commutative diagram
$$
\vcenter{
\xymatrix{
y \ar[r] \ar[d] & x_2 \ar[d] \\
x_1 \ar[r]      & x
}
}
\quad\text{lying over}\quad
\vcenter{
\xymatrix{
A_1 \times_A A_2 \ar[r] \ar[d] & A_2 \ar[d] \\
A_1 \ar[r]      & A.
}
}
$$
\item[(S2)]
The condition of (S1) holds for diagrams in $\mathcal{F}$ lying over
a diagram in $\mathcal{C}_\Lambda$ of the form
$$
\xymatrix{
          & k[\epsilon] \ar[d] \\
A  \ar[r] & k.
}
$$
Moreover, if we have two commutative diagrams in $\mathcal{F}$
$$
\vcenter{
\xymatrix{
y \ar[r]_c \ar[d]_a & x_\epsilon \ar[d]^e \\
x \ar[r]^d          & x_0
}
}
\quad\text{and}\quad
\vcenter{
\xymatrix{
y' \ar[r]_{c'} \ar[d]_{a'} & x_\epsilon \ar[d]^e \\
x \ar[r]^d                 & x_0
}
}
\quad\text{lying over}\quad
\vcenter{
\xymatrix{
A \times_k k[\epsilon] \ar[r] \ar[d] & k[\epsilon] \ar[d] \\
A  \ar[r] & k
}
}
$$
then there exists a morphism $b : y \to y'$ in
$\mathcal{F}(A \times_k k[\epsilon])$ such that $a = a' \circ b$.
\end{enumerate}
\end{definition}

\begin{definition}
\label{definition-minimal-versal}
Let $\mathcal{F}$ be a predeformation category.
We say a versal formal object $\xi$ of $\mathcal{F}$ is
{\it minimal}\footnote{This may be nonstandard terminology. Many
authors tie this notion in with properties of tangent spaces.
We will make the link in
Section \href{https://stacks.math.columbia.edu/tag/06IL}{section miniversal objects existence (Tag 06IL)}.}
if for any morphism of formal objects
$\xi' \to \xi$ the underlying map on rings is surjective.
Sometimes a minimal versal formal object is called {\it miniversal}.
\end{definition}

\begin{definition}
\label{definition-tangent-space}
Let $\mathcal{F}$ be a predeformation category.
The {\it tangent space $T \mathcal{F}$ of $\mathcal{F}$}
is the set $\overline{\mathcal{F}}(k[\epsilon])$
of isomorphism classes of objects in the fiber category $\mathcal
F(k[\epsilon])$.
\end{definition}

\noindent
The general notion of minimality introduced in
Definition \ref{definition-minimal-versal}
can sometimes be deduced from the behaviour on tangent spaces.
Let $\xi$ be a formal object of the predeformation category
$\mathcal{F}$ and let
$\underline{\xi} : \underline{R}|_{\mathcal{C}_\Lambda} \to \mathcal{F}$
be the corresponding morphism. Then we can consider the following
the condition
\begin{equation}
\label{equation-bijective}
d\underline{\xi} : \text{Der}_\Lambda(R, k) \to T\mathcal{F} 
\text{ is bijective}
\end{equation}
and the condition
\begin{equation}
\label{equation-bijective-orbits}
d\underline{\xi} : \text{Der}_\Lambda(R, k) \to T\mathcal{F} 
\text{ is bijective on }\text{Der}_\Lambda(k, k)\text{-orbits.}
\end{equation}
Here we are using the identification
$T\underline{R}|_{\mathcal{C}_\Lambda} = \text{Der}_\Lambda(R, k)$ of
Example \href{https://stacks.math.columbia.edu/tag/06IC}{example tangent space prorepresentable functor (Tag 06IC)}
and the action (\ref{equation-action}) of derivations
on the tangent spaces. If $k' \subset k$ is separable, then
$\text{Der}_\Lambda(k, k) = 0$ and the two conditions are equivalent.
It turns out that, in the presence of condition (S2) a versal formal
object is minimal if and only if $\underline{\xi}$ satisfies
(\ref{equation-bijective-orbits}).
Moreover, if $\underline{\xi}$ satisfies (\ref{equation-bijective}), then
$\mathcal{F}$ satisfies (S2).

\noindent
A special case of the constructions above are the map
\begin{equation}
\label{equation-map}
\gamma : \text{Der}_\Lambda(k, k) \longrightarrow T\mathcal{F}
\end{equation}
and the action
\begin{equation}
\label{equation-action}
a : \text{Der}_\Lambda(k, k) \times T\mathcal{F} \longrightarrow T\mathcal{F}
\end{equation}
defined for any predeformation category $\mathcal{F}$.
Note that if $\varphi : \mathcal{F} \to \mathcal{G}$ is a morphism
of predeformation categories, then we get commutative diagrams
$$
\vcenter{
\xymatrix{
\text{Der}_\Lambda(k, k) \ar[r]_-\gamma \ar[rd]_\gamma &
T\mathcal{F} \ar[d]_{d\varphi} \\
& T\mathcal{G}
}
}
\quad\text{and}\quad
\vcenter{
\xymatrix{
\text{Der}_\Lambda(k, k) \times T\mathcal{F} \ar[r]_-a
\ar[d]_{1 \times d\varphi} &
T\mathcal{F} \ar[d]^{d\varphi} \\
\text{Der}_\Lambda(k, k) \times T\mathcal{G} \ar[r]^-a &
T\mathcal{G}
}
}
$$

\noindent
The following observations are uninteresting in the classical case or when
$k/k'$ is a separable field extension, because then
$\text{Der}_\Lambda(k, k)$ and $\text{Der}_\Lambda(k, V)$ are zero.
There is a canonical identification
$$
\Mor_{\mathcal{C}_\Lambda}(k, k[\epsilon]) =
\text{Der}_\Lambda(k, k).
$$
Namely, for $D \in \text{Der}_\Lambda(k, k)$ let $f_D : k \to k[\epsilon]$
be the map $a \mapsto a + D(a)\epsilon$. More generally, given a finite
dimensional vector space $V$ over $k$ we have
$$
\Mor_{\mathcal{C}_\Lambda}(k, k[V]) =
\text{Der}_\Lambda(k, V)
$$
and we will use the same notation $f_D$ for the map associated to the
derivation $D$. We also have
$$
\Mor_{\mathcal{C}_\Lambda}(k[W], k[V]) =
\Hom_k(V, W) \oplus \text{Der}_\Lambda(k, V)
$$
where $(\varphi, D)$ corresponds to the map
$f_{\varphi, D} : a + w \mapsto a + \varphi(w) + D(a)$. We will sometimes write
$f_{1, D} : a + v \to a + v + D(a)$ for the automorphism
of $k[V]$ determined by the derivation $D : k \to V$. Note that
$f_{1, D} \circ f_{1, D'} = f_{1, D + D'}$.

\medskip\noindent
Let $\mathcal{F}$ be a predeformation category over $\mathcal{C}_\Lambda$.
Let $x_0 \in \Ob(\mathcal{F}(k))$. By the above there is a canonical
map
$$
\gamma_V :
\text{Der}_\Lambda(k, V)
\longrightarrow
\overline{\mathcal{F}}(k[V])
$$
defined by $D \mapsto f_{D, *}(x_0)$. Moreover, there is an action
$$
a_V : \text{Der}_\Lambda(k, V) \times \overline{\mathcal{F}}(k[V])
\longrightarrow
\overline{\mathcal{F}}(k[V])
$$
defined by $(D, x) \mapsto f_{1, D, *}(x)$. These two maps are compatible,
i.e., $f_{1, D, *}f_{D', *}x_0 = f_{D + D', *}x_0$ as follows from a
computation of the compositions of these maps. Note that the maps
$\gamma_V$ and $a_V$ are independent of the choice of $x_0$ as there
is a unique $x_0$ up to isomorphism.


\noindent
Explicitly, $p: \mathcal{F} \to \mathcal{C}$ is cofibered in groupoids if
the following two conditions hold:
\begin{enumerate}
\item For every morphism $f: U \to V$ in $\mathcal{C}$ and every object
$x$ lying over $U$, there is a morphism $x \to y$ of $\mathcal{F}$ lying
over $f$.
\item For every pair of morphisms $a: x \to y$ and $b: x \to z$
of $\mathcal{F}$ and any morphism $f: p(y) \to p(z)$ such that $p(b) = f
\circ p(a)$, there exists a unique morphism $c: y \to z$ of $\mathcal
F$ lying over $f$ such that $b = c \circ a$.
\end{enumerate}


\medskip\noindent
We denote by $k[\epsilon]$ the ring of dual numbers over $k$.  More
generally, for a $k$-vector space $V$, we denote by $k[V]$ the $k$-algebra
whose underlying vector space is $k \oplus V$ and whose multiplication is given
by $(a, v) \cdot (a', v') = (aa', av' + a'v)$.  When $V = k$, $k[V]$ is the ring
of dual numbers over $k$.  For any finite dimensional $k$-vector space $V$
the ring $k[V]$ is in $\mathcal{C}_\Lambda$.


\begin{lemma}
\label{lemma-construct-bijective-orbits}
Let $\mathcal{F}$ be a predeformation category satisfying
(S2) which has a versal formal object. Then its minimal versal
formal object satisfies (\ref{equation-bijective-orbits}).
\end{lemma}

\begin{proof}
Let $\xi$ be a minimal versal formal object for $\mathcal{F}$, see
Lemma \href{https://stacks.math.columbia.edu/tag/06T5}{lemma minimal versal (Tag 06T5)}.
Say $\xi$ lies over $R \in \Ob(\widehat{\mathcal{C}}_\Lambda)$.
In order to parse (\ref{equation-bijective-orbits}) we point out
that $T\mathcal{F}$ has a natural $k$-vector space structure
(see
Lemma \href{https://stacks.math.columbia.edu/tag/06IH}{lemma tangent space vector space (Tag 06IH)}),
that $d\underline{\xi} : \text{Der}_\Lambda(R, k) \to T\mathcal{F}$
is linear (see
Lemma \href{https://stacks.math.columbia.edu/tag/06IJ}{lemma k linear differential (Tag 06IJ)}),
and that the action of $\text{Der}_\Lambda(k, k)$ is
given by addition (see
Lemma \href{https://stacks.math.columbia.edu/tag/06SU}{lemma action linear (Tag 06SU)}).
Consider the diagram
$$
\xymatrix{
& \Hom_k(\mathfrak m_R/\mathfrak m_R^2, k) \\
K \ar[r] & \text{Der}_\Lambda(R, k) \ar[r]^{d\underline{\xi}} \ar[u] &
T\mathcal{F} \\
& \text{Der}_\Lambda(k, k) \ar[u] \ar[ru]
}
$$
The vector space $K$ is the kernel of $d\underline{\xi}$.
Note that the middle column is exact in the middle as it is dual to the
sequence (\href{https://stacks.math.columbia.edu/tag/06S6}{equation sequence (Tag 06S6)}). If (\ref{equation-bijective-orbits})
fails, then we can find a nonzero element $D \in K$ which
does not map to zero in $\Hom_k(\mathfrak m_R/\mathfrak m_R^2, k)$.
This means there exists an $t \in \mathfrak m_R$ such that
$D(t) = 1$. Set $R' = \{a \in R \mid D(a) = 0\}$. As $D$ is a derivation
this is a subring of $R$. Since $D(t) = 1$ we see that $R' \to k$
is surjective (compare with the proof of
Lemma \href{https://stacks.math.columbia.edu/tag/06H0}{lemma essential surjection (Tag 06H0)}).
Note that $\mathfrak m_{R'} = \Ker(D : \mathfrak m_R \to k)$
is an ideal of $R$ and $\mathfrak m_R^2 \subset \mathfrak m_{R'}$. Hence
$$
\mathfrak m_R/\mathfrak m_R^2 =
\mathfrak m_{R'}/\mathfrak m_R^2 + k\overline{t}
$$
which implies that the map
$$
R'/\mathfrak m_R^2 \times_k k[\epsilon] \to R/\mathfrak m_R^2
$$
sending $\epsilon$ to $\overline{t}$ is an isomorphism. In particular
there is a map $R/\mathfrak m_R^2 \to R'/\mathfrak m_R^2$.

\medskip\noindent
Let $\xi \to y$ be a morphism lying over $R \to R/\mathfrak m_R^2$.
Let $y \to x$ be a morphism lying over
$R/\mathfrak m_R^2 \to R'/\mathfrak m_R^2$.
Let $y \to x_\epsilon$ be a morphism lying over
$R/\mathfrak m_R^2 \to k[\epsilon]$. Let $x_0$ be the unique (up to unique
isomorphism) object of $\mathcal{F}$ over $k$.
By the axioms of a category cofibred in groupoids we obtain
a commutative diagram
$$
\vcenter{
\xymatrix{
y \ar[r] \ar[d] & x_\epsilon \ar[d] \\
x \ar[r]   & x_0
}
}
\quad\text{lying over}\quad
\vcenter{
\xymatrix{
R'/\mathfrak m_R^2 \times_k k[\epsilon] \ar[r] \ar[d] & k[\epsilon] \ar[d] \\
R'/\mathfrak m_R^2 \ar[r] & k.
}
}
$$
Because $D \in K$ we see that $x_\epsilon$ is isomorphic
to $0 \in \mathcal{F}(k[\epsilon])$, i.e., $x_\epsilon$ is
the pushforward of $x_0$ via $k \to k[\epsilon], a \mapsto a$.
Hence by
Lemma \href{https://stacks.math.columbia.edu/tag/06IS}{lemma lifting section (Tag 06IS)}
we see that there exists a morphism $x \to y$. Since
$\text{length}_\Lambda(R'/\mathfrak m_R^2) <
\text{length}_\Lambda(R/\mathfrak m_R^2)$
the corresponding ring map $R'/\mathfrak m_R^2 \to R/\mathfrak m_R^2$
is not surjective. This contradicts the minimality of
$\xi/R$, see part (1) of
Lemma \href{https://stacks.math.columbia.edu/tag/06T2}{lemma smallest where descends versal (Tag 06T2)}.
This contradiction shows that such a $D$ cannot exist, hence
we win.
\end{proof}
\end{document}
